\section{Related Work}
\label{sec:related_work}

\subsection{Image-Based Vulnerability Detection}

VulCNN~\cite{vulcnn} converts a function's PDG into a three-channel image using
degree, Katz, and closeness centrality scores, then classifies with a CNN.
It achieves scalability comparable to text-based approaches while retaining
program semantics, but the CNN cannot perform graph message-passing and
loses relational structure. VulGAI~\cite{vulgai} replaces Katz centrality with
second-order centrality and a jump-structured CNN architecture, improving speed
by $3.9\times$ over VulCNN. VulSCC~\cite{vulscc} augments image-based detection
with a frozen code LLM (Jina-embeddings-code) for semantic enrichment but
retains no graph branch. None of these methods incorporates graph message-passing.

\subsection{Graph-Based Vulnerability Detection}

ReGVD~\cite{regvd} builds a graph from unique code tokens or a sliding window
and applies a gated GNN initialized from a pre-trained language model.
iGnnVD~\cite{ignnvd} ensembles GCN, GAT, and APPNP over a bidirectional code
property graph. TensorGNN~\cite{tensorgnn} constructs a high-dimensional tensor
over AST, CFG, DFG, and NCS graphs processed by a tensor-based gated GNN.
All three are pure GNN approaches with no image or visual modality.

\subsection{Multimodal Vulnerability Detection}

TMF-Net~\cite{tmfnet} fuses contract graphs (GCN), opcode sequences (Bi-LSTM),
and code images (CNN) for smart-contract vulnerability detection. It is the
closest work to VulGCL in combining image and graph streams, but targets
Solidity smart contracts and cannot be directly applied to general C/C++/Java
source code.

\subsection{Stack Overflow and Code Security}

ICSD~\cite{icsd} detects insecure code snippets in Stack Overflow using a
Heterogeneous Information Network (HIN) and snippet2vec; it does not use a GNN
or image modality. iTrustSO~\cite{itrustso} extends ICSD with a GNN on the
snippet's parse graph integrated with HIN-based relations via hierarchical
attention (CodeHin2Vec); it classifies snippets posted on SO itself and uses
HIN structural relations, not community signals (votes, comments) as training
supervision. Chen et al.~\cite{chen2019reliable} demonstrate empirically that
SO vote scores and user reputation are unreliable security signals — insecure
answers average a score of 14 vs.\ 5 for secure answers — motivating the
signal-aware labeling strategy in VulGCL.

\subsection{Summary and Gap}

\begin{table}[t]
\centering
\caption{Comparison of related vulnerability detection methods.}
\label{tab:related}
\begin{tabular}{lcccc}
\toprule
Method & Graph & Image & LLM & General Code \\
\midrule
VulCNN~\cite{vulcnn}       & \texttimes & \checkmark & \texttimes & \checkmark \\
VulGAI~\cite{vulgai}       & \texttimes & \checkmark & \texttimes & \checkmark \\
VulSCC~\cite{vulscc}       & \texttimes & \checkmark & \checkmark & \checkmark \\
ReGVD~\cite{regvd}         & \checkmark & \texttimes & \texttimes & \checkmark \\
iGnnVD~\cite{ignnvd}       & \checkmark & \texttimes & \texttimes & \checkmark \\
TensorGNN~\cite{tensorgnn}  & \checkmark & \texttimes & \texttimes & \checkmark \\
TMF-Net~\cite{tmfnet}      & \checkmark & \checkmark & \texttimes & \texttimes \\
\midrule
\textbf{VulGCL (ours)}     & \checkmark & \checkmark & \checkmark & \checkmark \\
\bottomrule
\end{tabular}
\end{table}

As shown in Table~\ref{tab:related}, VulGCL is the first method to combine all
three modalities for general source-code vulnerability detection.
